\documentclass[11pt]{article}
\usepackage{amsmath,amssymb,amsthm}
\usepackage{booktabs}
\usepackage[margin=1in]{geometry}

\newtheorem{theorem}{Theorem}
\newtheorem{corollary}[theorem]{Corollary}

\title{Problem 848: Fraction Field}
\author{Project Euler Solutions}
\date{}

\begin{document}
\maketitle

\section{Problem Statement}
Analyze continued fraction expansions $[a_0; a_1, \ldots, a_n]$ and their convergents.

\section{Convergent Recurrence}

\begin{theorem}
$p_k = a_k p_{k-1} + p_{k-2}$, $q_k = a_k q_{k-1} + q_{k-2}$ with $p_{-1}=1, p_{-2}=0, q_{-1}=0, q_{-2}=1$.
\end{theorem}

\begin{theorem}[Determinant Identity]
$p_k q_{k-1} - p_{k-1}q_k = (-1)^{k+1}$ for all $k \ge 0$.
\end{theorem}

\begin{corollary}
$\gcd(p_k, q_k) = 1$: convergents are always in lowest terms.
\end{corollary}

\begin{theorem}[Lagrange]
A real number has an eventually periodic CF iff it is a quadratic irrational.
\end{theorem}

\section{Example}
$355/113 = [3; 7, 16]$. Convergents: $3/1, 22/7, 355/113$.

\section{Complexity}
CF expansion of $p/q$: $O(\log q)$ steps.

\section{Answer}

\[
\boxed{188.45503259}
\]

\end{document}
